\section{Related Work and Scientific Background}
    \subsection{Spiking Neural Networks}
    Spiking neural networks (SNNs) are computational models inspired by the communication and dynamics of biological neurons. Unlike conventional artificial neural networks (ANNs), which process continuous-valued signals, SNNs communicate through discrete events called spikes. Similar to biological neurons, spiking neurons have an internal state mostly represented by their membrane potential, that leaks over time and integrates incoming spikes. When the membrane potential crosses a certain threshold, the neuron generates an output spike, which is then transmitted to other neurons in the network. This sparse event-driven nature of SNNs allows them to process information in a more biologically plausible manner and can lead to energy-efficient computations, especially in neuromorphic hardware implementations.
    
    A defining characteristic of SNNs is that time is an explicit part of computation. Information can be represented by the number, timing or distribution of spikes across a population of neurons. The dynamics of membrane potentials and synaptic currents provide an intrinsic temporal memory allowing past inputs to influence subsequent responses.

    \subsubsection{LIF Spiking Neuron}
    The most commonly used neuron in SNNs is the Leaky Integrate-and-Fire (LIF) neuron. Much like the artificial neuron model, spiking neurons operate on a weighted sum of inputs, but rather than passing the result directly through a sigmoid or ReLU non-linearity, the weighted sum contributes to the membrane potential $U(t)$ of the neuron. If the neuron is sufficiently excited and reaches a threshold $\Theta$, the neuron will emit a spike onward and the membrane potential will discharge.

    The dynamics of the LIF can be modeled by an RC circuit, where the membrane potential $U(t)$ evolves according to the differential equation \cite{eshraghianTrainingSpikingNeural2023}:

    \begin{equation}
        \tau \frac{dU(t)}{dt} = -U(t) + I_{in}(t)R
        \label{eq:RC_circuit}
    \end{equation}

    where $\tau$ is the circuit time constant, $I_{in}(t)$ is the input current, and $R$ is the resistance. The solution to \eqref{eq:RC_circuit} is given by:

    \begin{equation}
        U(t) = I_{in}R + [U_0 - I_{in}R]e^{-\frac{t}{\tau}}
        \label{eq:RC_solution}
    \end{equation}

    Using the forward Euler method, the discrete-time solution can be approximated as:
    \begin{equation}
        U[t] = \beta U[t-1] + (1 - \beta) I_{in}[t] 
        \label{eq:RC_discrete}
    \end{equation}

    where $\beta = e^{-\frac{1}{\tau}}$ is the decay rate of the membrane potential $U[t]$.
    
    In deep learning, the weighting factor of an input is typically a learnable parameter. In SNNs, the input current $I_{in}(t)$ is often modeled as a weighted sum of incoming spikes from presynaptic neurons, so that $(1-\beta)I_{in}[t] = W X[t]$. We can then rewrite \eqref{eq:RC_discrete} as:
    
    \begin{equation}
        U[t] = \underbrace{\beta U[t-1]}_{\text{decay}} + \underbrace{W X[t]}_{\text{input}} - \underbrace{R[t-1]}_{\text{reset}}
        \label{eq:membrane_potential}
    \end{equation}

    If the neuron spiked at the previous time step, the reset term subtracts the threshold $\Theta$ from the membrane potential, effectively discharging it; otherwise, the reset term has no effect. A spike is generated if the membrane potential exceeds the threshold $\Theta$:

    \begin{equation}
        S_{out}[t] = 
        \begin{cases}
            1, & \text{if } U[t] \geq \Theta \\
            0, & \text{otherwise}
        \end{cases}
        \label{eq:spike_generation}
    \end{equation}


    \begin{figure}
        \centering
        \includegraphics[width=0.8\textwidth]{LIF_neuron.png}
        \caption{\textbf{Leaky Integrate-and-Fire (LIF) neuron driven by a constant input current.} The membrane potential $U[t]$ charges toward the threshold $\Theta$, emits a spike, and is reset by subtraction, producing a regular spike train.}
        \label{fig:LIF_neuron}
    \end{figure}
    

    \subsubsection{Spike Encodings}
    In an SNN, data is represented through \emph{temporal patterns of spikes}, rather than through continuous-valued activations. It is therefore necessary to encode the input data into a spike train. There are several input encoding schemes that can be used to encode the input information, each with its own advantages and trade-offs \cite{yamazakiSpikingNeuralNetworks2022,eshraghianTrainingSpikingNeural2023}:
    
    \begin{itemize}
        \item \textbf{Rate Coding:} The information is encoded in the average number of spikes generated by a neuron over a time window. For example, a bright pixel may generate a high firing rate, while a dark pixel generates a low firing rate. The firing rate can be deterministic or stochastic (Poisson process). Rate coding is considered more robust to noise and can be easily implemented, but it requires a longer time window to accurately represent the information.
        \item \textbf{Latency Coding:} The precise timing of spikes carries information. For example, with time-to-first-spike coding, a bright pixel produces an early spike, while a dark pixel spikes later or not at all. Latency coding assigns much more meaning to each individual spike. It is more energy efficient, but more sensitive to noise and require precise timing.
        \item \textbf{Population Coding:} Information is represented by the collective activity of a group of neurons over a time window, rather than the activity of individual neurons. In population coding, each neuron has its own unique receptive field, corresponding to its stimulus excitation over a range of continuous values. As a result, each neuron fires differently based on its receptive field. The information is encoded both in the pattern of which neurons in the population fired, and in how strongly each of them fired. Population coding has a richer representational capacity in a shorter time window and is more robust to noise, but it requires a larger number of neurons. Population encoding has been shown to be effective in control tasks \cite{tangDeepReinforcementLearning2020,huebotterSpikingNeuralNetworks2025a}.
    \end{itemize}

    \begin{figure}
    \centering
    \begin{subfigure}[b]{0.48\linewidth}
        \centering
        \includegraphics[width=\linewidth]{rate_encoding}
        \caption{}
        \label{fig:enc-rate}
    \end{subfigure}
    \hfill
    \begin{subfigure}[b]{0.48\linewidth}
        \centering
        \includegraphics[width=\linewidth]{latency_encoding}
        \caption{}
        \label{fig:enc-latency}
    \end{subfigure}
    \caption{\textbf{Spike encoding schemes.} (a) Rate coding: The information is encoded in the average number of spikes generated by a neuron over a time window. (b) Latency coding: The precise timing of spikes carries information.}
    \label{fig:encodings}
\end{figure}

\begin{figure}[!htbp]
    \centering
    \includegraphics[width=0.8\linewidth]{population_encoding}
    \caption{\textbf{Population coding: The information is represented by the collective activity of a group of neurons over a time window.} Each neuron has its own unique receptive field, corresponding to its stimulus excitation over a range of continuous values. Two symmetric values are encoded in the figure: $+0.3$ and $-0.3$. Although the center neuron (green neuron) emits the same spike rate for both values, the other excited neurons are different - preventing the ambiguity.}
    \label{fig:enc-population}
\end{figure}


\subsubsection{Training Spiking Neural Networks}
    Our goal is to optimize the learnable weights $W$ of the network with respect to some loss function $\mathcal{L}$. To achieve this, the gradient of the loss is computed with respect to each learnable parameter by applying the chain rule from the final layer back to each weight. The gradient is then used to update the weights such that the error decreases. If the gradient is zero, there is no weight update. Spikes are sparse, singular events in time, and are therefore non-differentiable. Although the internal membrane state of neurons keeps updating continuously, no gradient flows between the synapses of consecutive layers of the network. This is known as the \emph{dead neuron} problem in SNNs. 

    More formally, the spike function \eqref{eq:spike_generation} can be represented by the Heaviside step function $H$:
    \begin{equation}
        S[t] = H(U[t]-\Theta)
        \label{eq:heaviside-function}
    \end{equation}

    with its derivative
    \begin{equation}
        \frac{\partial S}{\partial U} = \delta(U-\Theta)
        \label{eq:heaviside-derivative}
    \end{equation}
    where $\delta$ is the Dirac delta function.

    Considering only the immediate influence of the input on the output, we get:
    \begin{equation}
        \frac{\partial\mathcal{L}}{\partial W}=\frac{\partial\mathcal{L}}{\partial S} \underbrace{\frac{\partial S}{\partial U}}_{\{ 0, \infty\}} \frac{\partial U}{\partial I} \frac{\partial I}{\partial W}
    \end{equation}

    From the recursive membrane potential \eqref{eq:membrane_potential} and the spike generation \eqref{eq:spike_generation}, SNNs are sequence-based models that resemble recurrent neural networks (RNNs) with implicit recurrence through $\beta$. Similar to RNNs, the network's computational graph can be unrolled in time. During training, a forward pass propagates the information through time to compute outputs and losses. During back propagation, the gradients of the unrolled graph are computed from the final error.

    

    
    More formally, the influence of the weights on the loss function is taken as the immediate influence, and the influence of all past


    Consider some loss function $\mathcal{L}$ and the network weights $W$; our goal is to find the derivative $\frac{\partial\mathcal{L}}{\partial W}$. The weights $W$ are applied at every time step $s$. As the data moves through the computational graph, each application of the weights at time $s$ will affect only present and future losses. The influence of all weight applications on present and future losses is summed together to define the global gradient \cite{eshraghianTrainingSpikingNeural2023}:

    \begin{equation}
        \frac{\partial\mathcal{L}}{\partial W}=\sum_{t}\frac{\partial\mathcal{L}[t]}{\partial W} = \sum_{t}\sum_{s\leq t}\frac{\partial \mathcal{L}[t]}{\partial W[s]} \frac{\partial W[s]}{\partial W}
        \label{eq:full-loss}
    \end{equation}

    As the weights are shared across all time steps: $W[0] = W[1] = W[2] = \dots = W$, $\partial W[s] / \partial W=1$, and \eqref{eq:full-loss} simplifies to:

    \begin{equation}
        \frac{\partial\mathcal{L}}{\partial W} = \sum_{t}\sum_{s\leq t}\frac{\partial \mathcal{L}[t]}{\partial W[s]}
        \label{eq:simplified-loss}
    \end{equation}


    Consider the immediate influence of the weights on the loss $\mathcal{L}$.
    As an example, let us backtrack one step in time to isolate the influence of a single step. That is:

    \begin{equation}
        \frac{\partial \mathcal{L}[t]}{\partial W[t-1]} = \frac{\partial \mathcal{L}[t]}{\partial S[t]}\frac{\partial S[t]}{\partial U[t]} \frac{\partial U[t]}{\partial U[t-1]} \frac{\partial U[t-1]}{\partial W[t-1]}
    \end{equation}

    The main obstacle in training SNNs with gradient-based methods is the spike generation function (Eq.~\ref{eq:spike_generation}). Being a Heaviside step function of the membrane potential, its derivative is zero everywhere except at the threshold $\Theta$, where it is undefined. As a result, the gradient of the loss with respect to the weights $W$ either vanishes or is undefined, and standard backpropagation cannot be applied directly (the \emph{dead neuron problem}). Three main families of approaches address this \cite{eshraghianTrainingSpikingNeural2023,yamazakiSpikingNeuralNetworks2022}. \textbf{ANN-to-SNN conversion} trains a conventional ANN and maps its activations to firing rates of an equivalent SNN; it achieves high accuracy but relies on rate coding and therefore requires many time steps. \textbf{Local learning rules}, such as spike-timing-dependent plasticity (STDP), update each weight based on the relative timing of pre- and post-synaptic spikes; they are biologically plausible and suited to on-chip learning, but scale poorly to deep networks and complex tasks. \textbf{Surrogate gradient learning} keeps the Heaviside function in the forward pass, but replaces its derivative in the backward pass with a smooth approximation, such as the derivative of a fast sigmoid, $\partial S_{out}/\partial U \approx 1/(1 + k|U - \Theta|)^2$, where $k$ controls the sharpness. Since the membrane potential (Eq.~\ref{eq:membrane_potential}) is recurrent in time, the network is unrolled over the simulation time steps and trained with backpropagation through time (BPTT), exactly as a recurrent neural network. Surrogate gradient learning allows SNNs to be trained end-to-end directly in the spiking domain, exploiting precise spike timing with few time steps, and is the approach adopted in this work.
The output layer of the a spiking

\subsubsection{Population Spiking Actor for Continuous Control}

\subsection{Drone Dynamics and Control}

\subsection{Theoretical Energy Consumption Calculation}
The theoretical inference energy of ANNs and SNNs can be approximated by counting the arithmetic operations each network performs and weighting every operation by its energy cost \cite{kimSpikingYOLOSpikingNeural2019,suDeepDirectlyTrainedSpiking2023}. We use the 45\,nm CMOS reference values of \cite{horowitz11ComputingsEnergy2014}, listed in Table~\ref{tab:op_energy}.

\begin{table}[htbp]
    \centering
    \begin{tabular}{lrr}
        \hline
        Operation & FLOPs & Energy \\
        \hline
        Addition (AC) & 1 & 0.9\,pJ \\
        Multiplication (MC) & 1 & 3.7\,pJ \\
        Multiply--accumulate (MAC) & 2 & 4.6\,pJ \\
        \hline
    \end{tabular}
    \caption{FLOP count and energy per FP32 operation in 45\,nm CMOS \cite{horowitz11ComputingsEnergy2014}. The MAC energy is the sum of the MC and AC costs.}
    \label{tab:op_energy}
\end{table}

\paragraph{ANN neuron.} A neuron $i$ with $N_{\mathrm{in}}$ inputs computes
\begin{equation}
    y_i = b_i + \sum_{j=1}^{N_{\mathrm{in}}} w_{ij} x_j ,
    \label{eq:ann_synops}
\end{equation}
which takes $N_{\mathrm{in}}$ multiplications, $N_{\mathrm{in}}-1$ additions to sum the products and one addition for the bias, i.e.\ $N_{\mathrm{in}}$ MACs. Its energy per inference is therefore
\begin{equation}
    E_{\mathrm{ANN}}^{\mathrm{neuron}} = N_{\mathrm{in}} E_{\mathrm{MAC}} .
    \label{eq:energy_ann_neuron}
\end{equation}

\paragraph{SNN neuron.} The inference cost is decomposed into spike-triggered synaptic operations and per-time-step neuron-state updates, while spike-generation and reset costs are neglected. At each time step, a LIF neuron therefore updates its membrane potential according to \eqref{eq:membrane_potential} without the reset term:
\begin{equation}
    U_i[t] = \beta U_i[t-1] + \sum_{j=1}^{N_{\mathrm{in}}} w_{ij} S_j[t] .
\end{equation}
Because the input spikes $S_j[t] \in \{0,1\}$ are binary, the synaptic sum needs no multiplications: each incoming spike adds its weight $w_{ij}$ to the potential with a single AC, and inputs that did not spike cost nothing. Let $fr_{\mathrm{in}}$ be the probability that an input spikes in a given step. Per time step, the neuron performs:
\begin{itemize}
    \item \textbf{Synaptic operations:} $fr_{\mathrm{in}} N_{\mathrm{in}}$ ACs on average.
    \item \textbf{Neuron-state update:} one MC for the leak $\beta U_i[t-1]$.
\end{itemize}
Over the $T$ time steps of one inference:
\begin{equation}
    E_{\mathrm{SNN}}^{\mathrm{neuron}} = T \left( fr_{\mathrm{in}} N_{\mathrm{in}} E_{\mathrm{AC}} + E_{\mathrm{MC}} \right).
    \label{eq:energy_snn_neuron}
\end{equation}

\paragraph{Limitations.} Equations~\eqref{eq:energy_ann_neuron} and \eqref{eq:energy_snn_neuron} give a theoretical approximation that assumes neuromorphic hardware, which processes only incoming spikes and skips inactive synapses; on a GPU the SNN runs densely and this saving is lost. 